\subsection{Выводы по главе}

Современные методы синтеза новых видов --- NeRF и 3DGS --- базируются на формулах прямого объёмного рендеринга, учитывающего поглощение и светимость среды. Так как изначально они создавались для снимков из реального мира, для которого эта модель слишком проста, поэтому вместо векторного поля $\mathcal (x, y, z) \to (\sigma_a, \sigma_e, p, \dots)$ методы предлагают использовать поля излучения $\mathcal (x, y, z, \theta, \phi) \to (R, G, B, \alpha)$, учитывающие направления взгляда.

В остальном оба этих метода основываются на восстановлении позиции камер и в случае 3DGS облака точек из COLMAP, рендеринге изображения, градиентном спуске по ошибке реконструкции по имеющимся изображениям. 